\scanpage{20}
% L1–5
Ce groupe s'identifie aussi au groupe des $f_0^*(\underline{\Lambda})$-automorphismes de $\underline{O}_X\mid U$, qui induisent l'identité dans $\underline{O}_X/\underline{I}.\underline{O}_X$ ; donc quand $U$ varie, ces groupes forment les sections d'un faisceau de groupes, noté
% L6
\[
\underline{\operatorname{Aut}}^{X_0}_T(X)\quad\text{ou}\quad\underline{\operatorname{Aut}}^{X_0}_{\underline{\Lambda}}(X)
\]
% L7–9
qui s'identifie aussi au \deleted{faisceau} de \deleted{\unsure{des}} groupes \deleted{des} \inserted{d'automorphismes de la} $f_0^*(\underline{\Lambda})$-\deleted{Algèbre} $\underline{O}_{X_0}$-augmentée $\underline{O}_X$.

% L10
D'autre part, considérons \deleted{pour tout ouvert} \inserted{le $S$-schéma}
% L11
\[
\deleted{\unsure{X}}\qquad\overline X=\underline{\operatorname{Hom}}_{T/S}(T\deleted{\unsure{_0}},X)\deleted{\unsure{,}}
\]
\margin{``schéma des sections de $X$ sur $\Lambda$''}
% L12–15
La \deleted{\unsure{bonne|homom.}} section $S\xrightarrow{\sim}T_0\longrightarrow T$ définit
\[
\underline{\operatorname{Hom}}_{T/S}(T,X)\longrightarrow\underline{\operatorname{Hom}}_{T_0/S}(T_0,X\times_T T_0)\simeq X_0,
\]
d'où un homomorphisme de $S$-\inserted{pré}schémas
\[
\overline X\xrightarrow{\pi_\Lambda^X}X_0
\]
% L16–18
Si $U$ est un ouvert de $X_{\deleted{\unsure{0}}\unsure{0|{}}}$, alors on a
\[
\deleted{U}\qquad\overline U\simeq(\pi_\Lambda^X)^{-1}(U)
\]
(compatible avec \unsure{la projection|les projections}).
% L19–20
Par suite, $\underline{\operatorname{Aut}}^{X_0}_T(X)$ \uline{opère comme faisceau d'automorphismes sur le} $X$-\uline{schéma} $\overline X$
% archive sheet number: 19

